% ============================================================
\part{Linux praktisch verstehen}
% ============================================================

% ── Kapitel 3 ────────────────────────────────────────────────
\chapter{Wir wollen eine Datei speichern und wiederfinden}

\kapitelziel{
Du lernst, wie Linux Dateien und Verzeichnisse organisiert und kannst
das Dateisystem vom Terminal aus sicher navigieren.
}
\kapitelstruktur

\section{Situation}
Du hast ein Python-Skript geschrieben und gespeichert. Jetzt öffnest
du ein neues Terminal -- und weißt nicht mehr, wo die Datei liegt.

\begin{aufgabeBox}
Navigiere im Dateisystem, finde deine Datei und lerne, Dateien
gezielt zu speichern und wiederzufinden.
\end{aufgabeBox}

\section{Erster Lösungsversuch}
Vielleicht kennst du Windows-Pfade wie \texttt{C:\textbackslash Users\textbackslash Wolfgang}.
Unter Linux sieht das anders aus.

\begin{problemBox}
Linux kennt keine Laufwerksbuchstaben. Alles beginnt mit \texttt{/}.
Wer das nicht weiß, versteht keine einzige Fehlermeldung.
\end{problemBox}

\begin{wissenBox}
Wir brauchen: den Aufbau des Linux-Dateisystems, die wichtigsten
Navigationsbefehle und das Konzept des aktuellen Verzeichnisses.
\end{wissenBox}

\section{Der pragmatische Weg}

\begin{lstlisting}[style=bash]
pwd                  # Wo bin ich gerade?
ls                   # Was ist hier?
ls -la               # Alles sehen, auch versteckte Dateien
cd /home/wolfgang    # Absoluter Pfad
cd ..                # Eine Ebene hoch
cd ~                 # Direkt zum Home-Verzeichnis
cd -                 # Zum letzten Verzeichnis zurueck
\end{lstlisting}

\begin{lstlisting}[style=text]
Linux-Dateisystem (vereinfacht):
/                    Root -- der Anfang von allem
├── home/
│   └── wolfgang/    Dein persoenliches Verzeichnis
├── etc/             Systemkonfiguration
├── var/             Logs und variable Daten
├── usr/             Installierte Programme
└── dev/             Geraetedateien (Sensoren, USB...)
\end{lstlisting}

\begin{loesungBox}
\textbf{Schritt 1:} Prüfe, wo du bist.
\begin{lstlisting}[style=bash]
pwd
# Ausgabe: /home/wolfgang
\end{lstlisting}
\textbf{Schritt 2:} Lege ein Projektverzeichnis an.
\begin{lstlisting}[style=bash]
mkdir -p ~/robotik/scripts
cd ~/robotik/scripts
\end{lstlisting}
\textbf{Schritt 3:} Erstelle eine Datei und prüfe sie.
\begin{lstlisting}[style=bash]
touch sensor_test.py
ls -la
\end{lstlisting}
\textbf{Schritt 4:} Dateien suchen.
\begin{lstlisting}[style=bash]
find ~ -name "sensor_test.py"
find ~ -name "*.py" -newer /tmp
\end{lstlisting}
\end{loesungBox}

\begin{merksatzBox}
\texttt{pwd} zeigt, wo du bist. \texttt{cd} bewegt dich. \texttt{ls} zeigt,
was da ist. Diese drei Befehle reichen für 80\% aller Navigation.
Die Tilde \texttt{\textasciitilde} ist immer dein Home-Verzeichnis.
\end{merksatzBox}

\begin{fehlerBox}
\begin{itemize}
  \item \textbf{Fehler:} \texttt{cd /home/wolfgang} auf einem fremden Rechner\\
        \textbf{Besser:} Immer \texttt{cd \textasciitilde} -- funktioniert überall
  \item \textbf{Fehler:} Leerzeichen im Verzeichnisnamen\\
        \textbf{Folge:} \texttt{cd mein projekt} scheitert -- Anführungszeichen nötig
\end{itemize}
\end{fehlerBox}

\begin{uebungBox}
\textbf{Aufgabe 1:} Navigiere in drei Schritten von \texttt{/} nach \texttt{/etc/default}.\\
\textbf{Aufgabe 2:} Erstelle die Struktur \texttt{\textasciitilde/robotik/band1/uebungen/} mit einem Befehl.\\
\textbf{Aufgabe 3:} Finde alle \texttt{.txt}-Dateien in \texttt{/etc}.
\end{uebungBox}

\begin{quizBox}
\begin{enumerate}
  \item Was gibt \texttt{pwd} aus?
  \item Was bedeutet \texttt{..} in einem Pfad?
  \item Wie lautet der absolute Pfad zum Home-Verzeichnis von Benutzer \texttt{ros}?
\end{enumerate}
\textbf{Antworten:} 1. Den absoluten Pfad des aktuellen Verzeichnisses.
2. Das übergeordnete Verzeichnis. 3. \texttt{/home/ros}
\end{quizBox}

\begin{haubeBox}
\textbf{Inodes}: Jede Datei hat einen Inode -- eine eindeutige Nummer, die
Metadaten speichert (Größe, Rechte, Zeitstempel). Der Dateiname ist nur
ein Zeiger auf den Inode. Daher können mehrere Namen auf dieselbe Datei
zeigen (\textit{Hardlinks}).

Das Gerät \texttt{/dev/ttyUSB0} ist eine Datei. Das Netzwerkinterface
\texttt{/proc/net/dev} ist eine Datei. Unter Linux ist \textit{alles} eine Datei.
\end{haubeBox}

\begin{robotikBox}
ROS2-Konfigurationsdateien liegen in \texttt{\textasciitilde/.ros/}, Logs in
\texttt{\textasciitilde/.ros/log/}. Wenn ein ROS2-Programm nicht startet,
sucht man zuerst dort. Das Dateisystem ist die Grundlage jeder Fehlersuche.
\end{robotikBox}

% ── Kapitel 4 ────────────────────────────────────────────────
\chapter{Wir wollen ein Programm starten}

\kapitelziel{
Du kannst Programme aus dem Terminal starten, Pfade verstehen und
Skripte ausführbar machen.
}
\kapitelstruktur

\section{Situation}
Du hast ein Python-Skript \texttt{roboter.py} gespeichert. Jetzt willst du
es starten -- aber \texttt{roboter.py} funktioniert nicht.

\begin{aufgabeBox}
Starte dein Python-Skript aus dem Terminal und verstehe, warum
\texttt{roboter.py} allein nicht ausreicht.
\end{aufgabeBox}

\begin{problemBox}
Das Terminal sucht Programme nur in bestimmten Verzeichnissen (\texttt{PATH}).
Ein Skript im aktuellen Verzeichnis wird nicht automatisch gefunden.
\end{problemBox}

\begin{wissenBox}
Wir brauchen: den PATH, den Shebang, Ausführungsrechte und den
Unterschied zwischen \texttt{python3 skript.py} und \texttt{./skript.py}.
\end{wissenBox}

\section{Der pragmatische Weg}

\begin{lstlisting}[style=bash]
# Methode 1: Interpreter direkt angeben (empfohlen)
python3 roboter.py

# Methode 2: Als ausführbare Datei (braucht Shebang + chmod)
chmod +x roboter.py
./roboter.py

# Wo sucht das Terminal nach Programmen?
echo $PATH
which python3
\end{lstlisting}

\begin{lstlisting}[style=python]
#!/usr/bin/env python3
# Erste Zeile: Shebang -- sagt Linux: dieses Skript mit python3 ausfuehren
print("Roboter gestartet")
\end{lstlisting}

\begin{loesungBox}
\begin{lstlisting}[style=bash]
# Shebang hinzufuegen (erste Zeile in roboter.py)
# Dann ausfuehrbar machen:
chmod +x roboter.py

# Starten -- Punkt-Slash sagt: "im aktuellen Verzeichnis suchen"
./roboter.py

# Oder ohne Shebang und chmod:
python3 roboter.py
\end{lstlisting}
\end{loesungBox}

\begin{merksatzBox}
\texttt{python3 skript.py} -- immer. \texttt{./skript.py} -- nur wenn Shebang
und \texttt{chmod +x} gesetzt. Das \texttt{./} ist kein Fehler, sondern Absicht:
\enquote{Suche hier, nicht im PATH.}
\end{merksatzBox}

\begin{fehlerBox}
\begin{itemize}
  \item \textbf{Fehler:} \texttt{roboter.py: command not found}\\
        \textbf{Grund:} Kein \texttt{./} und nicht im PATH
  \item \textbf{Fehler:} \texttt{Permission denied}\\
        \textbf{Grund:} \texttt{chmod +x} vergessen
\end{itemize}
\end{fehlerBox}

\begin{uebungBox}
\textbf{Aufgabe:} Erstelle ein Skript \texttt{hallo.py} mit Shebang und
\texttt{print("Hallo ROS2")}. Starte es auf beide Arten.
\end{uebungBox}

\begin{quizBox}
\begin{enumerate}
  \item Was speichert die Variable \texttt{\$PATH}?
  \item Warum braucht \texttt{./skript.py} ein \texttt{./}?
  \item Was macht \texttt{chmod +x datei}?
\end{enumerate}
\textbf{Antworten:} 1. Ordner, in denen Linux nach Programmen sucht.
2. Weil aktuelles Verzeichnis nicht im PATH ist.
3. Setzt das Ausführungsrecht.
\end{quizBox}

\begin{haubeBox}
\textbf{PATH im Detail:} \texttt{echo \$PATH} zeigt z.\,B.
\texttt{/usr/bin:/usr/local/bin:/home/ros/.local/bin}.
Linux durchsucht diese Ordner von links nach rechts. Wer eigene
Programme systemweit zugänglich machen will, kopiert sie nach
\texttt{/usr/local/bin} oder erweitert den PATH in \texttt{.bashrc}.
\end{haubeBox}

\begin{robotikBox}
ROS2-Befehle wie \texttt{ros2 run} funktionieren nur, wenn ROS2 im
PATH steht. Das geschieht durch:
\begin{lstlisting}[style=bash]
source /opt/ros/jazzy/setup.bash
\end{lstlisting}
Dieser Befehl fügt ROS2-Verzeichnisse zum PATH hinzu. In \texttt{.bashrc}
eingetragen, passiert das automatisch bei jedem Terminal-Start.
\end{robotikBox}

% ── Kapitel 5 ────────────────────────────────────────────────
\chapter{Wir wollen Dateien erstellen, kopieren, verschieben und löschen}

\kapitelziel{
Du kannst Dateien und Verzeichnisse mit Terminal-Befehlen verwalten
und kennst die wichtigsten Fallstricke beim Löschen.
}
\kapitelstruktur

\section{Situation}
Du willst eine Konfigurationsdatei sichern, bevor du sie änderst.
Außerdem müssen alte Logdateien gelöscht werden.

\begin{aufgabeBox}
Lege Kopien an, verschiebe Dateien in die richtige Struktur und
lösche gezielt -- ohne Datenverlust.
\end{aufgabeBox}

\begin{problemBox}
\texttt{rm} löscht sofort und endgültig. Es gibt keinen Papierkorb
im Terminal.
\end{problemBox}

\section{Der pragmatische Weg}

\begin{lstlisting}[style=bash]
# Erstellen
touch neue_datei.txt          # leere Datei
mkdir -p pfad/zu/ordner       # Ordner anlegen (p: auch Eltern)

# Kopieren
cp datei.txt backup.txt       # Datei kopieren
cp -r ordner/ ordner_backup/  # Ordner rekursiv kopieren

# Verschieben / Umbenennen
mv datei.txt neuer_name.txt   # Umbenennen
mv datei.txt ../anderer/pfad/ # Verschieben

# Löschen
rm datei.txt                  # Datei löschen
rm -r ordner/                 # Ordner mit Inhalt löschen
rm -i datei.txt               # Mit Rueckfrage (-i = interactive)

# Wildcards
cp *.py backup/               # Alle Python-Dateien kopieren
rm log_2024*.txt              # Alle alten Logs loeschen
\end{lstlisting}

\begin{loesungBox}
\textbf{Backup vor Änderung:}
\begin{lstlisting}[style=bash]
cp config.yaml config.yaml.bak
# Jetzt sicher bearbeiten
nano config.yaml
# Falls etwas schiefläuft:
cp config.yaml.bak config.yaml
\end{lstlisting}
\end{loesungBox}

\begin{merksatzBox}
Vor jedem \texttt{rm -r} einmal \texttt{ls} -- um sicher zu sein, was weg soll.
Ein \texttt{rm -rf /} löscht das gesamte System ohne Rückfrage. Linux
vertraut dir.
\end{merksatzBox}

\begin{fehlerBox}
\begin{itemize}
  \item \textbf{cp ordner ziel} statt \texttt{cp -r}: kopiert nur den Ordner, nicht den Inhalt
  \item \textbf{rm -rf \textasciitilde/}:  löscht dein komplettes Home-Verzeichnis
\end{itemize}
\end{fehlerBox}

\begin{uebungBox}
Erstelle folgende Struktur komplett mit Terminal-Befehlen:
\begin{lstlisting}[style=text]
~/robotik/
├── band1/
│   ├── uebungen/
│   └── loesungen/
└── backup/
\end{lstlisting}
Kopiere dann alle Dateien aus \texttt{uebungen/} nach \texttt{backup/}.
\end{uebungBox}

\begin{quizBox}
\begin{enumerate}
  \item Was ist der Unterschied zwischen \texttt{mv} und \texttt{cp}?
  \item Warum ist \texttt{rm -i} sicherer als \texttt{rm}?
  \item Was macht \texttt{mkdir -p a/b/c}, wenn \texttt{a} nicht existiert?
\end{enumerate}
\textbf{Antworten:} 1. mv verschiebt/benennt um; cp lässt Original bestehen.
2. Es fragt vor jedem Löschen nach. 3. Es erstellt alle drei Verzeichnisse.
\end{quizBox}

\begin{haubeBox}
\textbf{Hardlinks und Softlinks:} \texttt{ln datei link} erzeugt einen Hardlink --
zwei Namen für dieselben Daten. \texttt{ln -s datei link} erzeugt einen
Softlink (Symlink) -- ein Zeiger. ROS2 nutzt intensiv Symlinks, um
mehrere Workspace-Versionen zu verwalten.
\end{haubeBox}

\begin{robotikBox}
Rosbag-Dateien (\texttt{.db3}) können gigabytegross werden. Mit
\texttt{find \textasciitilde/.ros/log -mtime +7 -delete} löscht man automatisch
Logs, die älter als 7 Tage sind -- wichtig auf Robotern mit kleiner SSD.
\end{robotikBox}

% ── Kapitel 6 ────────────────────────────────────────────────
\chapter{Wir wollen wissen, was auf dem Computer gerade läuft}

\kapitelziel{
Du kannst laufende Prozesse anzeigen, identifizieren und beenden.
}
\kapitelstruktur

\section{Situation}
Ein ROS2-Node läuft noch im Hintergrund und blockiert einen Port.
Du weißt nicht, wie du ihn findest und stoppst.

\begin{aufgabeBox}
Finde alle laufenden Prozesse, identifiziere den störenden und
beende ihn gezielt.
\end{aufgabeBox}

\begin{problemBox}
Programme laufen nicht immer im Vordergrund. Hintergrundprozesse
sind unsichtbar -- bis man weiß, wo man schauen muss.
\end{problemBox}

\section{Der pragmatische Weg}

\begin{lstlisting}[style=bash]
ps aux                    # Alle Prozesse anzeigen
ps aux | grep python3     # Nur Python-Prozesse

top                       # Live-Übersicht (q zum Beenden)
htop                      # Komfortabler (apt install htop)

# Prozess beenden
kill 1234                 # PID 1234 beenden (SIGTERM)
kill -9 1234              # Sofort beenden (SIGKILL)
killall python3           # Alle python3-Prozesse beenden

# Prozesse im Hintergrund
python3 sensor.py &       # Im Hintergrund starten
jobs                      # Hintergrundprozesse anzeigen
fg %1                     # Prozess 1 in Vordergrund holen
\end{lstlisting}

\begin{loesungBox}
\begin{lstlisting}[style=bash]
# 1. Prozess finden
ps aux | grep ros2
# Ausgabe: ros  12345  ... python3 /opt/ros/jazzy/bin/ros2 run ...

# 2. PID merken (hier: 12345)
# 3. Beenden
kill 12345

# Falls er nicht reagiert:
kill -9 12345
\end{lstlisting}
\end{loesungBox}

\begin{merksatzBox}
Jeder Prozess hat eine eindeutige \textbf{PID} (Process ID). \texttt{kill}
sendet ein Signal -- standardmäßig SIGTERM (bitte beende dich). SIGKILL
(-9) ist der Notausschalter: sofortiger Abbruch ohne Aufräumen.
\end{merksatzBox}

\begin{fehlerBox}
\begin{itemize}
  \item \textbf{kill ohne Rechte:} Fremde Prozesse kann nur root beenden.
  \item \textbf{kill -9 zu früh:} Offene Dateien können korrupt werden.
\end{itemize}
\end{fehlerBox}

\begin{uebungBox}
Starte \texttt{python3 -c "import time; time.sleep(100)"} im Hintergrund.
Finde die PID und beende den Prozess.
\end{uebungBox}

\begin{quizBox}
\begin{enumerate}
  \item Was ist eine PID?
  \item Was ist der Unterschied zwischen SIGTERM und SIGKILL?
  \item Wie startet man ein Programm im Hintergrund?
\end{enumerate}
\textbf{Antworten:} 1. Process ID -- eindeutige Nummer eines laufenden Prozesses.
2. SIGTERM bittet zum Beenden; SIGKILL erzwingt es sofort.
3. Mit \texttt{\&} am Ende des Befehls.
\end{quizBox}

\begin{haubeBox}
\textbf{Prozess-Hierarchie:} Jeder Prozess hat einen Elternprozess.
\texttt{pstree} zeigt diese Baumstruktur. Wenn ein Elternprozess stirbt,
werden Kindprozesse oft zu \textit{Waisen} -- sie laufen weiter, bis
sie beendet werden. In ROS2 übernimmt der \texttt{ros2 launch}-Prozess
die Kontrolle über alle gestarteten Nodes.
\end{haubeBox}

\begin{robotikBox}
\texttt{ros2 node list} zeigt alle aktiven ROS2-Nodes. \texttt{ros2 node kill
/sensor\_node} beendet einen Node sauber. Bei hängendem System:
\begin{lstlisting}[style=bash]
pkill -f ros2
\end{lstlisting}
beendet alle Prozesse mit "ros2" im Befehlsnamen.
\end{robotikBox}

% ── Kapitel 7 ────────────────────────────────────────────────
\chapter{Wir wollen Fehlermeldungen verstehen}

\kapitelziel{
Du kannst Linux-Fehlermeldungen lesen, einordnen und die häufigsten
Probleme selbstständig lösen.
}
\kapitelstruktur

\section{Situation}
Du tippst einen Befehl -- und bekommst eine rote Fehlermeldung.
Was jetzt?

\begin{aufgabeBox}
Lerne, Fehlermeldungen zu lesen und systematisch zu beheben.
\end{aufgabeBox}

\begin{problemBox}
Fehlermeldungen sind präzise -- aber nur für jemanden, der das
Vokabular kennt. Ohne Grundwissen wirken sie einschüchternd.
\end{problemBox}

\begin{wissenBox}
Wir brauchen: Exit-Codes, die häufigsten Fehlertypen,
\texttt{stderr} vs. \texttt{stdout} und Log-Dateien.
\end{wissenBox}

\section{Die häufigsten Fehlermeldungen}

\begin{lstlisting}[style=bash]
# "command not found"
ros2 run my_robot start
# Ursache: ros2 nicht im PATH oder Tippfehler
# Loesung: source /opt/ros/jazzy/setup.bash

# "Permission denied"
./mein_skript.py
# Ursache: Kein Ausfuehrungsrecht
# Loesung: chmod +x mein_skript.py

# "No such file or directory"
cat /home/ros/config.yaml
# Ursache: Datei existiert nicht oder falscher Pfad
# Loesung: ls pruefen, Pfad korrigieren

# Exit-Code pruefen
python3 skript.py
echo $?    # 0 = Erfolg, alles andere = Fehler
\end{lstlisting}

\begin{loesungBox}
\textbf{Systematische Fehlersuche:}
\begin{lstlisting}[style=bash]
# 1. Fehlermeldung komplett lesen
# 2. Letzten Exit-Code pruefen
echo $?

# 3. Stderr von Stdout trennen
python3 skript.py 2>fehler.log
cat fehler.log

# 4. In Logs schauen
journalctl -xe          # Systemlogs
journalctl -f           # Live-Logs
tail -f ~/.ros/log/...  # ROS2-Logs
\end{lstlisting}
\end{loesungBox}

\begin{merksatzBox}
Exit-Code 0 bedeutet Erfolg. Jeder andere Wert ist ein Fehlercode.
\texttt{echo \$?} direkt nach einem Befehl zeigt seinen Exit-Code.
Fehlermeldungen kommen meist auf \texttt{stderr} -- dem \enquote{Fehlerkanal}.
\end{merksatzBox}

\begin{fehlerBox}
\begin{itemize}
  \item Fehlermeldung zu schnell wegklicken, ohne sie zu lesen.
  \item Nur die letzte Zeile lesen -- oft steht der echte Fehler weiter oben.
  \item Bei Python: der \texttt{Traceback} zeigt die genaue Zeile.
\end{itemize}
\end{fehlerBox}

\begin{uebungBox}
\textbf{Aufgabe:} Provoziere absichtlich drei verschiedene Fehler:
1. Datei, die nicht existiert, ausgeben.
2. Programm ohne Ausführungsrecht starten.
3. Python-Syntaxfehler. Lese und erkläre jede Fehlermeldung.
\end{uebungBox}

\begin{quizBox}
\begin{enumerate}
  \item Was bedeutet Exit-Code 0?
  \item Wohin gehen Fehlermeldungen normalerweise?
  \item Wie sehe ich den Exit-Code des letzten Befehls?
\end{enumerate}
\textbf{Antworten:} 1. Programm erfolgreich beendet.
2. Auf stderr (Standardfehlerausgabe). 3. Mit \texttt{echo \$?}
\end{quizBox}

\begin{haubeBox}
\textbf{Streams:} Jeder Prozess hat drei Standard-Streams:
\texttt{stdin} (Eingabe, 0), \texttt{stdout} (Ausgabe, 1), \texttt{stderr} (Fehler, 2).
Umleitung: \texttt{>} für stdout, \texttt{2>} für stderr, \texttt{\&>} für beides.
\begin{lstlisting}[style=bash]
python3 skript.py > ausgabe.txt 2>&1
\end{lstlisting}
Leitet alles (stdout + stderr) in eine Datei.
\end{haubeBox}

\begin{robotikBox}
ROS2-Nodes schreiben Fehler in strukturierte Log-Dateien unter
\texttt{\textasciitilde/.ros/log/}. Mit \texttt{ros2 bag info bag\_datei} lassen
sich auch aufgezeichnete Fehler rekonstruieren. Gute Fehlerkultur
beginnt damit, Meldungen komplett zu lesen statt sie zu ignorieren.
\end{robotikBox}
